% Cluster A: introduction
% Body-only LaTeX; preamble macros, cite keys, and cross-cluster labels per CONVENTIONS.md.

\section{Introduction}\label{sec:intro}

\subsection{The problem}\label{subsec:A-problem}

A recurring problem in complex dynamics is to connect a finite family of
dynamical continua while preserving selected accesses to their complements.
In renormalization theory the continua are small filled Julia sets or their
combinatorial stand-ins: the Douady--Hubbard theory of polynomial-like maps
produces them \cite{DH85}; the a~priori bounds of Kahn--Lyubich organize them
into decorated configurations whose limbs and accesses carry the analytic
estimates \cite{KL08}; Dudko's roadmap describes virtual Julia sets as
connected hulls of Cantor families of small Julia sets carrying a finite part
of the critical orbit \cite{Dud25}; and Kahn--Kapiamba--Lyubich work with a
Hubbard tree constructed inside the full dynamical carrier of a primitive
renormalization \cite{KKL26}.  In every one of these settings a proof needs,
at some stage, a single compact connected full continuum --- a
\emph{carrier} --- joining finitely many pairwise disjoint full continua
(\emph{atoms}, in the terminology of the source manuscript, fixed in
Section~\ref{sec:prelim}) inside a Jordan localization, avoiding a compact
forbidden set of obstacles and critical values, and, crucially, leaving
intact certain designated accesses: prime ends of the atoms through which
later curve-family, width, or landing arguments will read the picture.

Ordinary connectedness is far too weak a specification for this task,
because a connector chosen merely to connect can destroy the data that made
the connection worth having.  Three failure modes recur.  The connector, or
an excess set absorbed together with it, can \emph{absorb a protected side}:
the enlarged carrier swallows the collar through which a marked access was
to be read, and no curve in the new complement lands there from the marked
side any more.  It can \emph{identify two prime ends}: two designated
crosscut chains, distinct in the complement of the atom, can become
equivalent in the complement of the enlarged carrier, so that marked access
data collapses.  And it can \emph{create a separator}: a connector placed at
one stage can disconnect the free region needed at a later stage, and
monotone enlargement never undoes a separation.  All three failures are
invisible to the coarse bookkeeping --- homeomorphism types, degrees,
critical words --- that dynamical constructions usually carry, which is what
makes them dangerous.

\subsection{Main results}\label{subsec:A-results}

The first group of results makes the failures precise and shows that they
already occur in the simplest possible situations.
Theorem~\ref{thm:no-least} (manuscript Theorem~6.12) shows that there is no
least compact connected full continuum containing even the two-point package
$\{-1,1\}$: carrier selection by minimization is not merely inconvenient but
impossible, since minimal carriers abound while no least one exists.
Proposition~\ref{prop:access-collision} (manuscript Proposition~6.15)
exhibits a conformal, critical-value-free, degree-one branch of $z\mapsto
z^{2}$ in which the very first monotone absorption of pulled excess erases a
protected prime-end access, together with its conformally transported copy,
irreversibly.  Between them, the two obstructions close both naive routes: no
least hull exists to select, and a proposed hull cannot be blindly
completed.

The positive theory replaces hull selection by certification.  A
\emph{protected ribbon-routing certificate}
(Definition~\ref{def:certificate}; manuscript Definition~10.1) records, as
explicit data: the atoms; the frozen forbidden set; the protected accesses,
each presented as a prime-end crosscut chain with its selected sides,
truncated so that the terminal collars are pairwise disjoint; free-port arcs
attaching to each atom at a single free point; and an embedded ribbon
tree joining the ports, with its planar cyclic data declared.
Theorem~\ref{thm:realization} (manuscript Theorem~10.4) realizes every
certificate by an actual carrier: a compact connected full continuum,
compactly contained in the localization, containing the atoms, missing the
forbidden set, with singleton attachment traces at the atoms, in whose
complement every protected chain survives as a prime-end chain, the survival
assignment being injective and cyclic-order preserving.  No local
connectivity of any atom is assumed at any point: the hypothesis that such
constructions can never afford is never used.

The conceptual center of the paper is the exact criterion,
Theorem~\ref{thm:free-component} (manuscript Theorem~10.8).  Freeze, before
any connector is chosen, the complete finite future inventory --- atoms,
closed terminal protected collars, separated obstacles, critical-value
exclusions, and the localization-forbidden set --- call the union $Q$, and
let $\Omfree=U\setminus Q$ be the frozen free domain.  Then a protected
ribbon-routing certificate exists if and only if a single component of
$\Omfree$ admits simultaneous singleton ports to every required atom.
Moreover, the failure mode is monotone: if a first protected collar or
obstacle separates two required port classes, no further monotone connector
enlargement inside the localization can repair the failure.  The criterion
converts an existential search over embedded trees into a property of one
planar set fixed in advance, and it explains, in hindsight, the design of
the certificate: everything the future proof will read must be frozen into
$Q$ before routing begins, precisely so that the criterion quantifies over
nothing but $Q$.

The criterion is exact but not yet checkable from partial data, and the next
group of results addresses checkability.  Theorem~\ref{thm:common-domain}
(manuscript Theorem~10.11) gives a sufficient test in packet form: if the
frozen inventory can be assigned to finitely many pairwise disjoint full
continua --- augmented atoms and separated obstacles, each the filled hull of
a connected seed --- with one unprotected free port per original atom, then
the compacta form of Alexander duality forces the common complement to be
connected, producing the single free component that the exact criterion
demands.  Theorem~\ref{thm:audit} (manuscript Theorem~10.14) turns the
packet clauses into an audit for a fixed typed assignment: fullification is
independent of processing order; the clauses reduce to three verifiable
conditions (coverage of the frozen inventory by compact connected seeds;
pairwise disjoint completions compactly contained in the localization;
nonempty unprotected exposure locus for every original atom); and each
failure --- packet merger, localization escape, port loss --- is witnessed
by a certificate that persists under every monotone same-label enlargement.
Proposition~\ref{prop:cellwise} (manuscript Proposition~10.15) then gives
the most directly checkable sufficient condition: if the seeds sit in
pairwise disjoint closed Jordan cells, with one boundary-reaching
unprotected arc per atom, every packet clause follows.

The final substantive section applies the framework at full strength.  In
the direct-decoration setting --- a terminal primitive root continuum with
two disjoint univalent pullbacks generating a pullback compactification with
finite fibers and end fibers --- Theorem~\ref{thm:carrier} (manuscript
Theorem~10.45) produces, for any finite selection of decoration addresses
and finitely many protected root prime-end chains with proper total
impression, a protected regulated connector whose packet fullifications stay
inside pairwise disjoint gap cells, avoid every unselected finite fiber and
every end fiber, and preserve the transported prime ends together with their
cyclic order.  The application exercises every part of the machinery:
gap-cell isolation, packet replacement, impression ports, the audit, and the
realization theorem.

\subsection{What is new, and what is not}\label{subsec:A-honest}

The ingredients of the proofs are classical: the Janiszewski separation
theorem, Carath\'eodory's theory of prime ends, filled hulls, the Jordan and
Schoenflies extension theorems, and \v{C}ech cohomology with the compacta
form of Alexander duality.  Readers should not expect --- and will not
find --- a new theorem of classical planar topology.  The contribution is
the assembled certification framework: accesses protected explicitly, as
crosscut chains with selected sides, rather than implicitly, as intentions;
cyclic order carried as data and proved to survive; the complete future
inventory frozen before routing, so that existence becomes a property of one
planar set; typed packets whose completion is order-independent; and failure
certificates that are monotone-persistent, so that a negative answer is as
stable and as citable as a positive one.  In short, the paper contributes
proof technology for access-preserving planar connection, not a new fact
about the plane.  Section~\ref{sec:comparison} compares this technology in
detail with the standard continuum-hull constructions it is designed to
replace: least hulls, monotone absorption, regulated and Hubbard trees,
Moore-type monotone quotients, and the lamination-lift candidate.

\subsection{Relation to the source manuscript}\label{subsec:A-source}

All results are reproduced, with notational streamlining only, from the
author's research manuscript \cite{MWman}, and every reproduced statement
carries its manuscript number in its bracketed title.  There is exactly one
intentional mathematical divergence.  Manuscript Lemma~10.7, the slit lemma
underlying the exact criterion, is formally inconsistent as printed: its
slits are declared as maps into the open domain yet are required to take a
boundary value at one end, and the removal prescribed there also deletes the
free endpoints that the proof of manuscript Theorem~10.8 later joins inside
the complement.  Lemma~\ref{lem:slits} states and proves the repaired
form --- slits mapped into the closed domain, boundary-anchored, with only
the open slit interiors removed and the free tips retained for routing ---
and the proof of Theorem~\ref{thm:free-component} is rewritten accordingly.
Remark~\ref{rem:repair} records the divergence precisely.  No other
statement has been altered.  The comparison Section~\ref{sec:comparison} and
the worked examples of Sections~\ref{sec:criterion} and~\ref{sec:audit} are
new expository material for this paper; they promote no manuscript claim.

\subsection{Organization}\label{subsec:A-roadmap}

Section~\ref{sec:prelim} fixes notation and collects the classical
background: full continua and filled hulls, Janiszewski's theorem, prime
ends and crosscut chains with their selected sides, the compacta form of
Alexander duality, and the Jordan--Schoenflies extension package, together
with the manuscript's terms \emph{atom} and \emph{fullification}.
Section~\ref{sec:obstructions} proves the two motivating obstructions,
Theorem~\ref{thm:no-least} and Proposition~\ref{prop:access-collision}.
Section~\ref{sec:comparison} carries out the comparison with standard
continuum-hull constructions.  Section~\ref{sec:certificates} defines
protected ribbon-routing certificates and proves the realization theorem and
its conformal naturality.  Section~\ref{sec:criterion} proves the repaired
slit lemma and the exact free-component criterion, together with an example
separating individual access from the existence of a common free component.
Section~\ref{sec:packets} develops augmented-atom packets and the
common-domain theorem, and Section~\ref{sec:audit} proves order-independent
fullification, the port audit with its persistent failure certificates, and
the cellwise criterion, together with an example of a persistent
post-fullification packet merger.  Section~\ref{sec:decoration} develops the
finite direct-decoration application, culminating in
Theorem~\ref{thm:carrier}.  Section~\ref{sec:scope} states the scope
boundary.

\subsection{Scope}\label{subsec:A-scope}

A firm boundary is maintained throughout, following the source manuscript's
own discipline (Warning~10.59 of \cite{MWman}).  This paper works entirely
with finite, explicitly given planar data.  It does not construct a virtual
carrier or virtual component package; it does not construct or use a Wave
horizon; it establishes no analytic correspondence and no covering or
immersion action; it proves no uniform incidence bound and no compression;
and it makes no claim about local connectivity of the Mandelbrot set.  The
open problems of the source manuscript that would supply such objects --- in
particular manuscript Open Problems~10.6 and~10.58 --- are cited only as
open, and nothing here advances them.  What is proved is exactly this: once
a finite protected certificate or packet system is supplied, the planar
realization problem is completely solved, with an exact criterion, checkable
sufficient conditions, and persistent failure certificates.
Section~\ref{sec:scope} restates the boundary in detail.
